\documentclass{subfiles}
\begin{document}
    For the sake of simplicity the procedure for inserting a node in an RB tree is 
        almost the same to that of BST; the only difference is given by the 
        \lstinline{RB-Insert-Fixup} procedure at the end. See \Cref{RB-Insert}.

    % TODO: RB-Insert pseudo code 
    % \subfile{../../extras/TikZ/Procedure - RB -Insert}

    Consider the \lstinline{RB-Insert-Fixup} procedure shown below.
        By a simple analysis, six cases arise, in pairs symmetrical.
        Before we analize each, let us question which of the RB properties 
        may fail: is easely noted that property \ref{RB-trees:const1}, 
        \ref{RB-trees:const3} and \ref{RB-trees:const5} still hold.
        Thus, only property \ref{RB-trees:const2} and \ref{RB-trees:const4}
        can fail; the former failing when \(z\) is the root and the latter 
        when its parent is red. We consider cases 1-3 as 4-6 are symmetrical.

    \begin{description}
        \item [Case 1:] since both \(z.p\) and \(y\) are red;
            we can push the ``redness'' upwards towards \(z.p.p\).
            Since no other property fails, 
            we simply check if \(z.p.p\) break any property.

        \item [Case 2:] \(z\) is a right child while \(y\) is black;
            we can simply rotate on \(z.p\) and let Case 3 take care of the rest. 

        \item [Case 3:] \(z\) is a left child while \(y\) is black;
            whether executed directly or through Case 2, \(z.p.p\) identity is unchanged,
            then we can simply apply some color changes and a right rotation on \(z.p.p\).
    \end{description}

    For what regards \lstinline{RB-Insert} complexity: note that 
        it takes time proportional to the height of tree, 
        which we proved to be logarithmic; the \lstinline{RB-Insert-Fixup}
        is repeated only if case 1 gets executed. The total number of times 
        it can be executed is \(\log n\). Thus, \lstinline{RB-Insert} takes 
        \(\bigO{\log n}\) time.
\end{document}
